\chapter{Customer histories, heterogeneous effects and economic units}
\label{v131:histories-forests}

This chapter returns to the customer-level request after the network comparisons.
The earlier longitudinal, forest and dose studies supply the detailed methods;
the question here is what must remain unchanged when their outputs are reused.
Whole histories, within-period event order, the population being averaged and
the units of a treatment density can each change the economic answer. The
examples retain those distinctions without treating a common interface as a
common estimand.

\section{Resampling a sequence requires a sampling model}

Consider a bank evaluating a sequence of servicing decisions over twelve
months. The data contain repeated observations of each borrower, interventions
that depend on previous balances and competing events such as account closure.
The statistical unit might be a borrower history, a household, a branch or an
independent economic calendar. These choices imply different resampling laws.
A routine with ``bootstrap'' in its name does not determine the appropriate law.

LiNGAM's generic bootstrap resamples rows and refits the model. VARLiNGAM has
a distinct path: it fits a lag model, applies a contemporaneous LiNGAM model to
the residuals, resamples residual vectors and reconstructs a sequence using the
fitted autoregressive coefficients. The inspected loop temporarily sets the
lag-selection criterion to \texttt{None}. Consequently, one must inspect the
initial lag/coefficient settings and the loop's refit behavior before claiming
that it propagates model-selection uncertainty. \citep{v131_sw_lingam}

Residual resampling can be a sensible candidate under an appropriate innovation
law. It does not automatically preserve persistent residual dependence,
unmodeled common shocks or a different sampling calendar. Row resampling is a
different candidate suitable for a different observational unit. A banking
comparison should include a correctly specified innovation model and an adverse
case with persistent residuals, preserving the uncertainty target and failures
in both. Bootstrap edge frequencies should not be reported as a calibrated
effect confidence region without an additional argument.

pygformula's inspected bootstrap draws observed subject identifiers with
replacement, copies each subject's complete history and assigns a fresh replicate
identifier to each draw. It updates histories and refits covariate and outcome
models. Expanding to the requested simulation population occurs in a separate
step. These operations distinguish the observed sampling unit from the Monte
Carlo population used to evaluate a fitted regime. \citep{v131_sw_pygformula}

The fresh identifiers matter. Drawing borrower 17 twice should produce two
bootstrap histories, not concatenate them into one impossible person with two
records at every date. An integration test needs duplicate draws, unequal
history lengths and multiple interventions. It must also preserve the original
calendar when that calendar is part of the estimand. Subject resampling alone
does not account for a market shock shared by every subject on the same date.

Failed fits require equally explicit handling. The inspected bootstrap helper
catches an exception, warns and returns a dictionary whose
\texttt{boot\_results} is \texttt{None}. The caller filters such entries when
assembling one result list, but subsequently constructs \texttt{bootests} by
indexing every original entry. A failed entry therefore exposes a static
inconsistency in that caller path. This observation requires reproduction in
the pinned runtime before attributing an observed frequency or impact to it.
The integration must count attempted, successful and failed replications and
preserve the failure reason, irrespective of whether the eventual repair
aborts, returns an explicit unavailable interval or uses a preregistered rule.
\citep{v13_sw_pygformula,v131_sw_pygformula}

\section{The intervention clock and competing-event target}

The inspected simulation path applies an intervention to the current simulated
frame and then refreshes history variables before outcome prediction. Static
schedules index their values by the absolute time index \texttt{int\_values[t]}.
Threshold interventions bound the current action between the supplied limits.
Custom interventions receive the current frame, available pool and time index.
The API permits useful regimes, but does not itself certify that a custom
function uses only information available at the decision time.
\citep{v131_sw_pygformula}

For example, ``contact once delinquency has persisted for two months'' refers
to the observed history before the action, whereas ``contact borrowers who
will be delinquent next month'' uses future information. Both can be expressed
as executable code if future fields have been made available. A valid interface
must enforce the declared information clock and treat prohibited reads as a
contract violation. Thresholds should also retain their monetary units and
whether they change an action, a measurement or an eligibility condition.

Competing events alter the target even when the same outcome model is used.
In a stipulated discrete-time model, let $d_t$ be the probability of a competing
event in period $t$ among those still at risk, and let $y_t$ be the outcome
hazard conditional on avoiding that competing event. If competing events occur
before outcomes within a period, the cumulative incidence through $T$ is
\[
 F_Y(T)=\sum_{t=1}^T
 \left\{\prod_{s<t}(1-d_s)(1-y_s)\right\}(1-d_t)y_t.
\]
This formula follows by partitioning on the period of the first outcome event.
It differs from a hypothetical law that eliminates the competing event.

For two periods with $y=(0.1,0.2)$ and $d=(0.05,0.1)$, the first contribution
is $0.095$ and the second is $0.9(0.95)(0.9)(0.2)=0.1539$. Total incidence is
$0.2489$. Removing competing events while holding these outcome hazards fixed
gives $1-(0.9)(0.8)=0.28$. At a stipulated 10,000 accounts and CNY 100 per event,
the corresponding expected-loss difference is CNY 31,100. This is exact
illustrative arithmetic. It does not identify the effect of eliminating a
real bank's competing event or establish that the outcome hazard would remain
fixed under that intervention. \citep{v131_project_semantics}

The inspected pygformula source separately computes survival risks with and
without a competing-event path, using subject-level cumulative products and
then population means. That code provides a concrete place to check the
within-period ordering, history alignment and target definition. It does not
settle whether observed censoring or competing-event assumptions are valid in
the proposed banking application. \citep{v131_sw_pygformula}

\section{Following an R forest operation into its C++ implementation}

For grf, the causal-forest binding assigns the same variable index to treatment
and instrument, then calls an instrumental predictor. The predictor factory
constructs \texttt{InstrumentalPredictionStrategy}. The native prediction uses
a ratio of centered outcome--instrument and treatment--instrument moments.
When instrument and treatment coincide, this corresponds to a weighted local
covariance divided by a weighted local treatment variance. This traced route
supports the interpretation of continuous-treatment output as a local slope
under its functional assumptions. \citep{v131_sw_grf}

It does not identify an arbitrary finite change in a nonlinear dose-response
curve. A change from an offer rate of one unit to two units needs the appropriate
counterfactual contrast and exposure support. A slope, an average slope and a
finite contrast can coincide in a linear model and diverge otherwise. The bank
must choose the target before comparing software outputs.

The native variance routine forms influence-like quantities across tree
groups, separates between-group variation from group noise and calls a Bayes
debiaser. These groups concern the forest's computation. They are distinct
from the observed borrower clusters used by the R average-effect calculation.
Inspecting both paths avoids the false inference that a tree-group variance
option automatically accounts for a persistent banking calendar. The debiaser,
forest sampling, serialization and numerical safeguards require qualification
before a direct native interface is admitted. Section~\ref{exec13:grf-adapter}
develops the selected native representation, variance-field and persistence
contracts; source inspection does not replace their runtime tests. \citep{v131_sw_grf}

The maintained R interface is therefore a reasonable first integration
candidate, subject to a complete license/dependency and behavioral review.
A direct C++ port would have to preserve the forest representation, prediction
strategy, variance method, randomization and exceptional outcomes across version
changes. Avoiding one language boundary does not by itself justify that
maintenance commitment. The open-source decision should compare the operational
cost of an R process or library bridge with the full scientific and technical
cost of a maintained port.

\section{Population averaging, support and unavailable combinations}

In grf's inspected all-population AIPW path, the mean is a weighted average of
\texttt{get\_scores} output. The variance calculation sums centered weighted
scores within the declared clusters, squares these group contributions and
applies a group-count correction. For binary treatment, the score uses the
usual propensity residual divided by $\widehat e(1-\widehat e)$. For continuous
treatment, an auxiliary regression forest estimates the conditional variance
of the treatment; the debiasing weight divides the treatment residual by that
estimate. \citep{v131_sw_grf}

Small conditional variance is an information problem. Borrowers with essentially
one available offer rate do not support precise estimation of a local rate
response simply because a forest can fit the outcome. Qualification should
retain weak-support cases and identify whether the operation returns a
warning, an extreme weight, an unavailable result or a misleading finite value.
The inspected formula alone supplies no finite-sample guarantee for such cases.

The all, treated, control and overlap populations must remain distinct. In the
overlap path the code fits a residualized outcome on a residualized treatment
and uses cluster or heteroskedastic covariance as applicable. Moving to an
overlap target may improve stability while changing which banking population
the result describes. A wrapper must not silently make that choice on the
analyst's behalf. \citep{v131_sw_grf}

Option support must be determined from reachable control flow. The inspected
average-effect routine rejects TMLE when clusters or sample weights are supplied
before reaching later formulas that contain cluster-related expressions. Seeing
those later formulas would not justify advertising clustered TMLE at this entry
point. Likewise, only specified combinations support continuous treatment.
The adapter's capability matrix needs explicit accepted combinations and
typed refusal for the others. \citep{v131_sw_grf}

\section{A generalized propensity score has physical units}

Let a continuous exposure follow a location-scale law
$A=\mu(X)+s(X)R$, with $s(X)>0$. Conditional change of variables gives
\[
 f_{A\mid X}(a\mid x)
 =\frac{1}{s(x)}f_{R\mid X}
       \left(\frac{a-\mu(x)}{s(x)}\middle|x\right).
\]
The Jacobian $1/s(x)$ is necessary for a density with respect to the original
exposure measure. If a common residual density is used instead of
$f_{R\mid X}$, an additional residual-distribution assumption is being made.

CausalGPS's inspected \texttt{compute\_resid} returns the standardized residual
$(a-b)/c$. Its \texttt{compute\_density} helper calls R's density estimate and
interpolates it with \texttt{rule=2}. In the inspected kernel branch of
\texttt{estimate\_gps}, the standardized-residual density is assigned to the
GPS field without division by the predicted scale in that branch. The normal
branch instead evaluates a normal density using its fitted mean and scale.
These observations establish a concrete scale question for downstream matching,
weighting and dose-response construction. Section~\ref{sec:exec24-dose-response-reuse}
traces the selected weighting ratio and outcome frontend. Its scale witness
narrows the outstanding question without qualifying every downstream operation. \citep{v13_sw_causalgps,v131_sw_causalgps}

At a standardized residual of zero, a standard normal density is about
$0.3989423$. Two strata with scales one and two have original-exposure densities
$0.3989423$ and $0.1994711$ at their respective means. Retaining only the common
standardized-residual value loses that difference. If the exposure is measured
in units one hundred times smaller, its numerical value is multiplied by one
hundred and the corresponding density must be divided by one hundred. The
probability of any physical exposure interval remains unchanged. These exact
transformations provide a reference for a future unit-consistency test.
\citep{v131_project_semantics}

Some downstream procedures might deliberately use a transformed balancing
quantity, while others require a density in the original units. The integration decision therefore remains conditional on a declared use of
the returned field and a justified interpretation of its scale. The selected
weighting path has been traced; its remaining scale and native-conformance
questions must be resolved before admission. Silently changing the helper
could alter package behavior without resolving the scientific contract. A
matched comparison should examine heterogeneous scales, near-zero scales,
counterfactual dose evaluation and an explicit unit conversion, retaining any
unsupported cases.

The causaldrf comparison has a separate source obligation. Its accompanying
license notice identifies the year and copyright holder, supplementing the
package's license declaration. That resolves the availability of that notice;
it does not settle transitive licensing or qualify its interval routines.
Section~\ref{sec:exec24-dose-response-reuse} develops the selected interval-index
and pointwise-versus-simultaneous restrictions, which remain part of the
proposed comparison. Source availability, permission to distribute and justified
uncertainty remain separate questions. \citep{v131_sw_causaldrf}

\section{Implications for reusable causal operations}

These examples support a common semantic contract with explicit graph
orientation, query/failure type, independent observational unit, decision clock,
target population, dose units and uncertainty construction. A business-facing
assistant can elicit these fields and explain unresolved choices, while a
deterministic validator checks the declared operation against its supported
combinations. An LLM's plausible description is not evidence that a package
uses the intended covariance or intervention target.

The relevant next qualification compares complete causal requests. It retains
the same estimand and information across packages, separates sampling from
Monte Carlo error, includes failed fits and preserves unbounded sets. Numerical
agreement in a benign case is useful evidence about implementation, but cannot
establish assignment truth, identification or useful precision in adverse
banking regimes. Those remain scientific and institutional requirements.


The commercial comparison that follows applies the same criteria to documented
vendor capabilities and to this proposal's actual evidence. A product feature
is useful only for a specified request; a proposed preserving interface is not
yet demonstrated superiority over an existing engine.
